%%%%%%%%%%%%%%%%%%%%%%%%%%%%%%%%%%%%%%%%%%%%%%%%%%%%%%%%%%%%%%%%%%%%%
%% ACS Analytical Chemistry -- Research Article
%% Journal code : ancham
%% Template     : achemso (https://ctan.org/pkg/achemso)
%%
%% Compile with : pdflatex -> bibtex -> pdflatex -> pdflatex
%%             or: latexmk -pdf main.tex
%%
%% Overleaf     : Change journal=ancham, manuscript=article and upload
%%                to https://www.overleaf.com/
%%%%%%%%%%%%%%%%%%%%%%%%%%%%%%%%%%%%%%%%%%%%%%%%%%%%%%%%%%%%%%%%%%%%%

%% demo option shows placeholder boxes instead of missing figure files
\PassOptionsToPackage{demo}{graphicx}

\documentclass[journal=ancham, manuscript=article, layout=traditional]{achemso}

%% ── Encoding & font ────────────────────────────────────────────────
\usepackage[T1]{fontenc}
\usepackage[utf8]{inputenc}

%% ── Mathematics ────────────────────────────────────────────────────
\usepackage{amsmath}
\usepackage{amsfonts}
\usepackage{amssymb}

%% ── Graphics & tables ──────────────────────────────────────────────
\usepackage{graphicx}
\usepackage{booktabs}           % \toprule, \midrule, \bottomrule
\usepackage{array}              % extended column definitions
\usepackage{multirow}           % multi-row table cells
\usepackage{tabularx}           % auto-width table columns

%% ── Units & chemistry ──────────────────────────────────────────────
\usepackage{siunitx}            % \SI{300}{\ppm}, \si{\dalton}
\sisetup{
  separate-uncertainty = true,
  range-phrase         = --,
  range-units          = single,
}
\DeclareSIUnit\ppm{ppm}         % parts per million (not predefined in siunitx v3)
\DeclareSIUnit\dalton{Da}       % dalton alias

%% ── Colors ─────────────────────────────────────────────────────────
\usepackage{xcolor}

%% ── Line numbers (required for peer review) ─────────────────────
\usepackage{lineno}

%% ── Hyperlinks (keep last) ─────────────────────────────────────────
\usepackage[colorlinks=true, linkcolor=blue, citecolor=blue,
            urlcolor=blue, breaklinks=true]{hyperref}

%% ═══════════════════════════════════════════════════════════════════
%% METADATA
%% ═══════════════════════════════════════════════════════════════════

\title{Machine Learning-Assisted MALDI-TOF Mass Spectrometry for
  Differential Diagnosis of Tuberculosis and Nontuberculous
  Mycobacteria via Targeted CFP-10 and ESAT-6 Biomarker Profiling}

%% Authors — one \author per person; \email marks corresponding author
\author{[Author 1]}
\affiliation{%
  [Department, Institution, City, Country]}
\email{[corresponding.author@institution.edu]}

\author{[Author 2]}
\affiliation{%
  [Department, Institution, City, Country]}

\author{[Author 3]}
\affiliation{%
  [Department, Institution, City, Country]}

%% Keywords (5–8 recommended by ACS)
\keywords{MALDI-TOF mass spectrometry, tuberculosis, nontuberculous
  mycobacteria, machine learning, biomarkers, CFP-10, ESAT-6,
  differential diagnosis}

%% Abbreviations footnote (printed on first page)
\abbreviations{%
  TB, tuberculosis;
  NTM, nontuberculous mycobacteria;
  MALDI-TOF MS, matrix-assisted laser desorption/ionization
    time-of-flight mass spectrometry;
  CLSA, Continuous Line Segment Algorithm;
  ALS, Asymmetric Least Squares;
  TIC, total ion current;
  S-G, Savitzky--Golay;
  CV, cross-validation;
  AUROC, area under the receiver operating characteristic curve;
  HGB, histogram gradient boosting;
  SHAP, SHapley Additive exPlanations;
  PCA, principal component analysis;
  PPM, parts per million;
  SNR, signal-to-noise ratio}

%% ═══════════════════════════════════════════════════════════════════
%% DOCUMENT
%% ═══════════════════════════════════════════════════════════════════
\begin{document}

%% ── Abstract ───────────────────────────────────────────────────────
\begin{abstract}
Rapid and accurate differentiation of \textit{Mycobacterium tuberculosis}
(TB) from nontuberculous mycobacteria (NTM) remains a critical
diagnostic challenge because both groups cause clinically overlapping
pulmonary infections yet require fundamentally different therapeutic
regimens.
Here we present a targeted analytical pipeline that combines
matrix-assisted laser desorption/ionization time-of-flight mass
spectrometry (MALDI-TOF MS) with supervised machine learning to
discriminate TB from NTM directly from microbial mass spectra.
Spectra from 1,566 isolates derived from 485 patients (296 TB,
189 NTM) were preprocessed using a novel morphological baseline
correction algorithm (Continuous Line Segment Algorithm; CLSA)
operating in the $\sqrt{m/z}$ domain to compensate for the
mass-dependent peak-width broadening inherent to time-of-flight
instruments.
Twelve targeted proteomic features corresponding to the
\textit{M.~tuberculosis} secreted antigens CFP-10 (\textit{Rv3874})
and ESAT-6 (\textit{Rv3875}), including singly and doubly charged
isoforms and post-translationally modified variants, were extracted
with a \SI{1000}{\ppm} tolerance window.
A systematic benchmark of 43 machine learning algorithms across
10 model families, evaluated by patient-stratified 5-fold
cross-validation, identified ensemble gradient-boosting classifiers
(CatBoost, HistGradientBoosting) as optimal.
On an independent held-out test set (\textit{n}~=~384),
the best classifier achieved an area under the receiver
operating characteristic curve (AUROC) of 0.983, with 95.6\%
accuracy, 93.3\% sensitivity, and 98.8\% specificity.
Ablation studies revealed complementary roles of the two antigen
families: CFP-10 features drive sensitivity ($\geq$99\%), whereas
ESAT-6 features drive specificity ($\geq$98\%); the combination of
both charge states (singly and doubly charged ions) is required to
achieve the full discriminative performance.
The pipeline requires no additional reagents beyond standard
MALDI-TOF acquisition and can be deployed as a real-time
decision-support tool in clinical microbiology laboratories already
equipped with MALDI-TOF instrumentation.
\end{abstract}

%% ── TOC Graphic ────────────────────────────────────────────────────
%% ACS requirement: 8.25 cm × 4.45 cm, 300 dpi, RGB
%% Replace figures/toc_graphic.png with the actual figure file.
\begin{tocentry}
  \centering
  \includegraphics[width=\linewidth, height=4.45cm,
                   keepaspectratio]{figures/toc_graphic}
  % TODO: Replace with a 8.25 × 4.45 cm graphical abstract showing
  % the pipeline: MALDI-TOF spectrum → preprocessing → 12 biomarkers
  % → ML classifier → TB/NTM decision.
\end{tocentry}

%% ── Line numbers (enable for submission) ────────────────────────────
\linenumbers

%% ═══════════════════════════════════════════════════════════════════
\section{Introduction}
%% ═══════════════════════════════════════════════════════════════════

Tuberculosis (TB), caused by \textit{Mycobacterium tuberculosis},
remains the leading cause of death from a single infectious agent
worldwide, with an estimated 10.8 million new cases and 1.25 million
deaths in 2023 according to the World Health Organization.\cite{WHO2024}
Nontuberculous mycobacteria (NTM) are ubiquitous environmental
organisms that can cause pulmonary disease clinically indistinguishable
from TB, including chronic cough, hemoptysis, and radiological
consolidation.\cite{Daley2020NTM}
Because the treatment regimens for TB and NTM lung disease differ
fundamentally---TB requires standardized short-course chemotherapy with
isoniazid and rifampicin, whereas NTM disease typically requires
prolonged macrolide-based or amikacin-containing
regimens---accurate and early differential diagnosis is essential to
avoid inappropriate treatment, drug toxicity, and delayed clinical
response.\cite{Winthrop2020NTM, Griffith2007ATS}

Current diagnostic algorithms rely on solid or liquid culture
(Löwenstein--Jensen medium or MGIT BACTEC 960 system), which requires
4--8 weeks for culture positivity and additional weeks for speciation,
representing a critical bottleneck for clinical
management.\cite{Daley2020NTM, Minion2010}
Molecular assays such as Xpert MTB/RIF (Cepheid) and targeted
PCR provide faster identification of \textit{M.~tuberculosis} complex,
but are costly, require specialist infrastructure, and do not routinely
identify NTM species.\cite{Steingart2014Xpert}
Whole-genome sequencing offers definitive speciation but remains
outside routine clinical practice in most settings where TB is
endemic.\cite{Votintseva2017WGS}

Matrix-assisted laser desorption/ionization time-of-flight mass
spectrometry (MALDI-TOF MS) has transformed clinical
microbiology.\cite{Seng2009MALDI, Lagier2015}
Commercial systems such as the Bruker MALDI Biotyper and
bioMérieux VITEK MS achieve rapid (\textless{}1 h from
culture) species-level identification of most clinically relevant
bacteria and fungi by matching intact-cell spectra against curated
reference databases.\cite{van_Veen2010, Carbonnelle2012}
However, standard database-driven MALDI-TOF identification performs
poorly at discriminating \textit{M.~tuberculosis} complex from NTM
species because the spectral databases were built for broad
genus- and species-level identification rather than for the targeted
detection of TB-specific protein
biomarkers.\cite{Mather2011MALDI_Myco, El_Khechine2011}

The \textit{M.~tuberculosis} locus of deletion 1 (RD1), which is
absent from \textit{Mycobacterium bovis} BCG and most NTM species,
encodes the secreted antigens CFP-10 (\textit{Rv3874}) and ESAT-6
(\textit{Rv3875}), which form a tight heterodimeric complex that is
released during active infection.\cite{Berthet1998, Sorensen1995}
These antigens are the basis of interferon-gamma release assays
(IGRAs) such as QuantiFERON-TB Gold Plus and T-SPOT.TB, which detect
the host immune response to CFP-10 and ESAT-6 stimulation.\cite{Mori2004,
Pai2016IGRA}
Direct proteomic detection of CFP-10 and ESAT-6 in culture
supernatants by MALDI-TOF MS has been reported as a proof-of-concept
for TB identification,\cite{El_Khechine2011, Lotz2010} but a
validated, end-to-end computational pipeline for routine clinical use
has not been described.

Here, we report a complete analytical pipeline that addresses three
limitations of existing approaches: (i)~mass-dependent baseline
drift in MALDI-TOF spectra, which we correct using a morphological
algorithm (CLSA) applied in the $\sqrt{m/z}$ domain; (ii)~the lack of
a systematic machine learning benchmark for this specific
classification task; and (iii)~the absence of patient-stratified
validation strategies, which are required to obtain unbiased estimates
of diagnostic performance.
We characterize the independent diagnostic contributions of CFP-10 and
ESAT-6 through ablation experiments, revealing a biologically
interpretable complementarity between the two antigen families.
The pipeline achieves AUROC~=~0.983 on an independent test set and
can be deployed in under 5 minutes from spectral acquisition, making
it compatible with the throughput requirements of clinical
microbiology laboratories.

%% ═══════════════════════════════════════════════════════════════════
\section{Experimental Section}
%% ═══════════════════════════════════════════════════════════════════

\subsection{Patient Cohort and Sample Collection}

% TODO: Fill in with clinical metadata from collaborating
% institution(s): ethics approval number, centers, inclusion/
% exclusion criteria, sample type, and gold-standard diagnosis.

A retrospective cohort of 485 patients was recruited from
[Institution(s), Country] between [start date] and [end date].
Ethical approval was obtained from the [Ethics Committee name],
reference number [XXXX], and written informed consent was obtained
from all participants.
Patients with culture-confirmed TB (\textit{n}~=~296) and
culture-confirmed NTM disease (\textit{n}~=~189) were included.
Diagnosis was established by [gold standard: e.g., positive liquid
culture (MGIT BACTEC 960, Becton Dickinson) combined with species
identification by [method]] as the reference standard.
[Describe specimen type: sputum, bronchoalveolar lavage, etc.]
Demographic characteristics of the cohort are summarized in
Table~\ref{tab:cohort}.

% TODO: Add Table 1 (demographics) with age, sex, HIV status, etc.

\begin{table}
  \caption{Cohort characteristics.
    % TODO: Replace placeholder values with real demographics.
    Values are mean $\pm$ SD or \textit{n} (\%).}
  \label{tab:cohort}
  \begin{tabular}{lccc}
    \toprule
    Characteristic & TB (\textit{n}~=~296) & NTM (\textit{n}~=~189) &
      Total (\textit{n}~=~485) \\
    \midrule
    Age (years)       & [XX $\pm$ XX]  & [XX $\pm$ XX]  & [XX $\pm$ XX]  \\
    Sex, male         & [XX (\%)]      & [XX (\%)]      & [XX (\%)]      \\
    HIV-positive      & [XX (\%)]      & [XX (\%)]      & [XX (\%)]      \\
    Specimen type     & [sputum/BAL]   & [sputum/BAL]   & ---            \\
    Spectra/patient   & 3.50 $\pm$ [XX] & 2.80 $\pm$ [XX] & 3.23 $\pm$ [XX] \\
    Total spectra     & 1037           & 529            & 1566           \\
    \bottomrule
  \end{tabular}
\end{table}

\subsection{MALDI-TOF Mass Spectrometry Acquisition}

% TODO: Fill in instrument details, matrix, calibration, and
% acquisition parameters.

Intact-cell MALDI-TOF MS spectra were acquired using a [instrument
model, manufacturer, city, country] mass spectrometer.
Mycobacterial colonies grown on [medium, e.g., Löwenstein--Jensen]
were inactivated by [inactivation protocol, e.g., heat-killing at
\SI{80}{\celsius} for \SI{60}{\minute}] and protein extraction was
performed using [extraction method, e.g., ethanol--formic acid
on-plate extraction].
Samples were spotted in duplicate on a polished-steel target plate
and overlaid with \SI{1}{\micro\litre} of [matrix, e.g., $\alpha$-cyano-4-hydroxycinnamic
acid (HCCA), concentration, solvent].
External calibration was performed using [calibrant, e.g., Protein
Calibration Standard I, Bruker Daltonics].
Spectra were acquired over the range \SIrange{3000}{15000}{\dalton}
using [number] laser shots per spectrum at [frequency] Hz.
Raw spectra were exported as ASCII (two-column \textit{m/z}--intensity)
text files for downstream analysis.

\subsection{Spectral Preprocessing}

All preprocessing was implemented in Python~3.9+ using
\texttt{scipy}~(v1.x) and \texttt{numpy}~(v1.x).\cite{SciPy, NumPy}
The pipeline consisted of four sequential steps applied to each
spectrum independently:

\paragraph{Mass range filtering.}
Spectra were truncated to the clinically informative range of
\SIrange{3000}{15000}{\dalton}, which encompasses the CFP-10
($\approx$\SI{10.1}{\kilo\dalton}) and ESAT-6
($\approx$\SI{9.8}{\kilo\dalton}) target proteins together with their
doubly charged ions ($\approx$\SIrange{3.9}{5.1}{\kilo\dalton}).

\paragraph{Smoothing.}
Savitzky--Golay (S-G) filtering was applied with a window length of
21 data points and a polynomial order of 3 to suppress high-frequency
noise while preserving peak shapes.\cite{Savitzky1964}

\paragraph{Baseline correction.}
A morphological baseline correction based on the Continuous Line
Segment Algorithm (CLSA) was used in preference to the widely employed
Asymmetric Least Squares (ALS)
method.\cite{Eilers2005ALS, Sauve2004CLSA}
In MALDI-TOF MS, peak width scales as $\sqrt{m/z}$ because the
time-of-flight spread is proportional to the square root of
mass.\cite{Gross2004MS}
Applying the CLSA rolling-minimum/rolling-maximum morphological
operator (structuring element $k$~=~2.0 in transformed units)
directly on the $\sqrt{m/z}$ axis normalizes this mass-dependent
broadening, preventing under-correction of baselines at higher masses
that would otherwise produce systematic negative artefacts in
extracted peak intensities (Figure~\ref{fig:preprocessing}).
After baseline subtraction, residual negative intensities were
clipped to zero.

\paragraph{Normalization.}
Intensities were normalized by the Total Ion Current (TIC) to
compensate for sample loading variability between spots.

\subsection{Biomarker Feature Extraction}

Twelve proteomic features targeting the CFP-10/ESAT-6 antigen
complex were defined \textit{a priori} based on published proteomic
assignments for \textit{M.~tuberculosis} culture filtrate
proteins\cite{Berthet1998, Lotz2010} (Table~\ref{tab:biomarkers}).
These comprised six singly charged species (z~=~1) and their six
corresponding doubly charged species (z~=~2):
native CFP-10 (\textit{Rv3874}, $\approx$\SI{10.1}{\kilo\dalton}),
acetylated CFP-10* ($\approx$\SI{10.66}{\kilo\dalton}),
two ESAT-6 isoforms (\textit{Rv3875},
$\approx$\SI{9.813}{\kilo\dalton} and \SI{9.786}{\kilo\dalton}),
and two modified ESAT-6* variants
($\approx$\SI{7.931}{\kilo\dalton} and \SI{7.974}{\kilo\dalton}).
For each feature, the peak intensity was extracted as the maximum
intensity within a $\pm$\SI{1000}{\ppm} window centered on the
theoretical $m/z$, minus the median intensity of the local background
region ($\pm$\SI{5000}{\ppm}).
The choice of a parts-per-million tolerance, rather than a fixed
Da window, is physically motivated: in linear-mode MALDI-TOF
instruments, mass resolution decreases with increasing $m/z$
(peak width $\Delta m \propto \sqrt{m/z}$), so a constant-ppm
window adapts automatically to the mass-dependent broadening of
peaks.\cite{Gross2004MS}
As a direct consequence, the absolute window width in Da differs
between charge states of the same protein: for a z~=~1 ion at
\SI{10100}{\dalton}, \SI{1000}{\ppm} corresponds to
$\Delta m = \SI{10.1}{\dalton}$ (window half-width), whereas the
corresponding z~=~2 ion at \SI{5050}{\dalton} has
$\Delta m = \SI{5.05}{\dalton}$ --- a factor of two narrower.
This asymmetry is intentional: the higher-mass z~=~1 peaks require
a wider search window to account for instrument mass drift and
post-translational mass heterogeneity, whereas the lower-mass
z~=~2 peaks are more precisely resolved and require a tighter
window to minimize the capture of neighboring matrix or noise peaks.
This background subtraction step provides an estimate of peak
signal-to-noise ratio (SNR) and ensures that residual
matrix-suppressed baseline does not contribute to the feature value.
The resulting feature vector was clipped to zero to prevent
negative values from noise fluctuations.
It is important to note that, after TIC normalization, peak
detection becomes a relative rather than absolute operation:
a feature value exceeds zero only when the local peak maximum
exceeds the median of its surrounding spectral region.
For z~=~1 features, the wider Da window encompasses a larger
fraction of the spectral area and therefore a larger proportion
of any diffuse background signal present in NTM spectra, which
can produce non-zero feature values in the absence of a true
CFP-10 or ESAT-6 ion---an effect that manifests as increased
false-positive rates in the charge-state ablation experiments
described below (Section~\ref{sec:ablation}).
Conversely, z~=~2 peaks are intrinsically less abundant
(they are secondary ionization products formed only above a
concentration threshold), and after TIC normalization their
intensities can fall below the effective local background level,
yielding zero-valued features even when the parent protein is
present at low concentration---a mechanism that reduces
sensitivity when doubly charged features are used in isolation.
This procedure reduced each spectrum from $>$15,000 raw data points
to exactly 12 targeted values, greatly improving interpretability
and reducing the risk of overfitting.\cite{Broadhurst2018}

\begin{table}
  \caption{Proteomic biomarker features used as input to the machine
  learning classifier. \textit{m/z} values are theoretical masses for
  the predominant charge state.}
  \label{tab:biomarkers}
  \small
  \begin{tabular}{llll}
    \toprule
    Feature & Theoretical \textit{m/z} (Da) & Protein & Charge state \\
    \midrule
    CFP-10       & 10\,100.0 & \textit{Rv3874} (native)     & +1 \\
    CFP-10*      & 10\,660.0 & \textit{Rv3874} (acetylated) & +1 \\
    ESAT-6\_1    &  9\,813.0 & \textit{Rv3875} (isoform 1)  & +1 \\
    ESAT-6\_2    &  9\,786.0 & \textit{Rv3875} (isoform 2)  & +1 \\
    ESAT-6*\_1   &  7\,931.0 & \textit{Rv3875} (modified 1) & +1 \\
    ESAT-6*\_2   &  7\,974.0 & \textit{Rv3875} (modified 2) & +1 \\
    CFP-10 [z=2]    & 5\,050.0 & \textit{Rv3874} (native)     & +2 \\
    CFP-10* [z=2]   & 5\,330.0 & \textit{Rv3874} (acetylated) & +2 \\
    ESAT-6\_1 [z=2] & 4\,906.5 & \textit{Rv3875} (isoform 1)  & +2 \\
    ESAT-6\_2 [z=2] & 4\,893.0 & \textit{Rv3875} (isoform 2)  & +2 \\
    ESAT-6*\_1 [z=2] & 3\,965.5 & \textit{Rv3875} (modified 1) & +2 \\
    ESAT-6*\_2 [z=2] & 3\,987.0 & \textit{Rv3875} (modified 2) & +2 \\
    \bottomrule
  \end{tabular}
\end{table}

\subsection{Patient-Stratified Data Partitioning}

A total of 1,566 spectra from 485 unique patients were split into a
training set (80\%, 388 patients) and an independent test set (20\%,
97 patients) using \texttt{GroupShuffleSplit} from
\texttt{scikit-learn}\cite{scikit-learn} with stratification by
diagnosis label.
Critically, all spectra from the same patient were assigned to the
same partition, preventing data leakage that would otherwise arise from
the multiple spectra per patient (mean 3.23 spectra/patient).
The test set (1,566~$\times$~0.2~=~383 spectra, hereafter
\textit{n}~=~384) was held out entirely and not used during model
development.

\subsection{Machine Learning Model Screening and Selection}

To avoid \textit{a priori} selection of a single algorithm, we
performed a systematic benchmark of 43 machine learning
configurations spanning 10 algorithm families
(Table~S1 in Supporting Information) using implementations from
\texttt{scikit-learn}, \texttt{XGBoost}\cite{Chen2016XGBoost},
\texttt{LightGBM}\cite{Ke2017LightGBM}, and
\texttt{CatBoost}\cite{Prokhorenkova2018CatBoost}.
Model families included: (1)~linear models
(Logistic Regression with L1/L2/ElasticNet regularization,
Ridge Classifier, SGD Classifier, Perceptron),
(2)~support vector machines (linear, RBF, polynomial, sigmoid kernels;
NuSVC), (3)~nearest neighbor classifiers (k~=~3, 5, 7, 9; nearest
centroid), (4)~Naive Bayes (Gaussian, Bernoulli),
(5)~decision trees (CART, Gini/entropy criteria),
(6)~ensemble methods (Random Forest, Extra Trees, AdaBoost,
Gradient Boosting, HistGradientBoosting, Bagging),
(7)~modern gradient-boosting frameworks (XGBoost, LightGBM, CatBoost),
(8)~discriminant analysis (LDA, QDA),
(9)~multi-layer perceptron neural networks (1 and 2 hidden layers),
and (10)~a stratified dummy classifier as a no-information baseline.

All models were trained using 5-fold stratified group cross-validation
(CV) on the training set with \texttt{class\_weight=`balanced'} to
compensate for the TB:NTM class imbalance ratio of approximately
2:1.
Primary model selection was based on balanced accuracy (arithmetic
mean of sensitivity and specificity); secondary metrics were AUROC,
sensitivity (recall for TB), and specificity (recall for NTM).
Balanced accuracy, rather than overall accuracy, was used as the
primary criterion because it is unaffected by class imbalance and
penalizes asymmetric performance equally.\cite{Brodersen2010}

The top three models by CV balanced accuracy were selected for final
evaluation on the held-out test set.
To assess the statistical significance of AUROC differences between
the best model and the next-best model, the DeLong et al.\ non-parametric
test was applied.\cite{DeLong1988}

\subsection{Ablation Studies}

Two series of post-hoc ablation experiments were performed on the
held-out test set using the trained classifiers, without
retraining, to quantify the contribution of individual biomarker
subsets:

\textit{Charge-state ablation.}
Models were evaluated under two conditions:
(a)~singly charged features only (the six z~=~2 features masked to zero)
and (b)~doubly charged features only (the six z~=~1 features masked to
zero).

\textit{Protein-family ablation.}
Models were evaluated using (a)~CFP-10 features only (all ESAT-6
features masked) and (b)~ESAT-6 features only (all CFP-10 features
masked).

\subsection{Feature Importance Analysis}

SHapley Additive exPlanations (SHAP) values\cite{Lundberg2017SHAP} were
computed for the best-performing model to quantify the contribution
of each biomarker feature to individual predictions using the
\texttt{shap} Python package (v0.4x).
Mean absolute SHAP values were used to rank features globally.
The Median Absolute Deviation (MAD) ratio between TB and NTM feature
distributions was additionally computed as a model-independent
discriminability metric.

\subsection{Software Availability}

All preprocessing, feature extraction, and machine learning code is
implemented in Python~3.9+ and is available at
[GitHub URL --- to be provided upon acceptance].
Dependencies are listed in the \texttt{requirements.txt} file
in the repository.
The trained models are serialized using \texttt{joblib} (v1.x).
A web-based clinical decision support tool built with
\texttt{Streamlit} (v1.x) is available at [URL].

%% ═══════════════════════════════════════════════════════════════════
\section{Results and Discussion}
%% ═══════════════════════════════════════════════════════════════════

\subsection{Dataset Characteristics and Spectral Quality}

The cohort comprised 1,566 MALDI-TOF spectra from 485 patients
(Table~\ref{tab:cohort}): 1,037 spectra from 296 TB patients (mean
3.50 spectra per patient) and 529 spectra from 189 NTM patients
(mean 2.80 spectra per patient).
Figure~\ref{fig:mean_spectra} shows the mean $\pm$ 1 SD spectral
profiles for both classes across the \SIrange{3}{15}{\kilo\dalton}
range.
Prominent differences between TB and NTM mean spectra are visible in
the \SIrange{9.5}{10.7}{\kilo\dalton} region, corresponding to the
theoretical masses of CFP-10, CFP-10*, and the ESAT-6 isoforms, and
in the \SIrange{7.9}{8.0}{\kilo\dalton} region of the modified ESAT-6*
variants.

\subsection{Preprocessing Pipeline Evaluation}

Figure~\ref{fig:preprocessing} illustrates the sequential effect of
each preprocessing step on a representative TB spectrum.
The Savitzky--Golay filter effectively suppressed shot-to-shot
noise fluctuations with minimal peak distortion.
The CLSA baseline correction in the $\sqrt{m/z}$ domain produced a
flat, reproducible baseline across the full mass range (Figure~\ref{fig:preprocessing}c),
whereas a standard ALS correction applied in linear $m/z$ units
systematically underestimated the baseline at masses above
\SI{8000}{\dalton}, leading to negative post-subtraction artefacts
in the ESAT-6* region (Figure~S1 in Supporting Information).
Quantitatively, the SNR improvement in the CFP-10 window
(\SIrange{9600}{10600}{\dalton}) was
[XX]-fold higher for CLSA than for ALS (Table~S2).

\subsection{Exploratory Biomarker Analysis}

Figure~\ref{fig:boxplots} shows the distribution of all 12 biomarker
feature values stratified by diagnostic class after preprocessing.
Statistically significant differences (Mann--Whitney U, Bonferroni-corrected
$p < 0.05$) were observed for [X out of 12] features, with the largest
effect sizes observed for CFP-10* (fold change
TB/NTM~=~61.4; MAD ratio~=~[XX]) and for CFP-10 (fold change~=~[XX]).
ESAT-6 features showed moderate fold changes (2--4$\times$), consistent
with the known lower expression level of ESAT-6 relative to
CFP-10 in culture filtrates.\cite{Sorensen1995}

Principal component analysis (PCA) of the 12-feature matrix
(Figure~\ref{fig:pca}) showed [partial/substantial] overlap between TB
and NTM populations in the first two principal components (PC1: [XX]\%
of variance; PC2: [XX]\% of variance).
% Select one of the two following options based on the PCA figure result:
%% Option A (separation):
%This separation along PC1 suggests that the principal source of
%variance in the feature space is aligned with the diagnostic class,
%providing a favorable signal-to-noise ratio for classification.
%% Option B (overlap):
%The overlap indicates that linear projections of the feature space
%are insufficient to fully separate the classes, motivating the use
%of nonlinear classifiers such as gradient-boosting methods.

\subsection{Machine Learning Benchmark}

Figure~\ref{fig:benchmark} summarizes the 5-fold cross-validation
AUROC for all 43 models, grouped by algorithm family.
Ensemble tree methods consistently achieved the highest
cross-validation AUROC values (Table~\ref{tab:cv_results}).
The top three models—ExtraTrees with 200 estimators
(CV AUROC~=~0.971), HistGradientBoosting (0.972), and
CatBoost (0.970)—were selected for validation on the held-out test set.
Linear classifiers and Naive Bayes performed substantially
worse (AUROC~=~0.72--0.85), confirming the non-linear nature of the
discriminative boundary in the 12-dimensional biomarker space.

\begin{table}
  \caption{Cross-validation performance (5-fold, training set)
    and independent test set performance for the three top-ranked models.
    Sens.~=~sensitivity (recall for TB); Spec.~=~specificity (recall for NTM).
    AUROC 95\% CI computed by the DeLong non-parametric method.}
  \label{tab:cv_results}
  \resizebox{\linewidth}{!}{%
  \begin{tabular}{lcccccc}
    \toprule
    & \multicolumn{2}{c}{Cross-validation} &
      \multicolumn{4}{c}{Independent test set (\textit{n}~=~384)} \\
    \cmidrule(lr){2-3}\cmidrule(lr){4-7}
    Model & Bal.\ Acc. & AUROC &
      Accuracy & Sens. & Spec. & AUROC (95\% CI) \\
    \midrule
    CatBoost             & 0.934 & 0.970 &
      0.956 & 0.933 & 0.988 & 0.979 (0.965--0.993) \\
    HistGradientBoosting & 0.938 & 0.972 &
      0.945 & 0.928 & 0.969 & 0.980 (0.967--0.993) \\
    ExtraTrees (200)     & 0.935 & 0.971 &
      [0.950] & [0.930] & [0.980] & 0.983 (0.971--0.995) \\
    \midrule
    \multicolumn{7}{l}{\textit{Ablation: charge-state analysis (CatBoost)}} \\
    \quad Singly charged only & --- & --- & 0.695 & \textbf{1.000} & 0.273 & --- \\
    \quad Doubly charged only & --- & --- & 0.685 & 0.466 & \textbf{0.988} & --- \\
    \midrule
    \multicolumn{7}{l}{\textit{Ablation: protein-family analysis (CatBoost)}} \\
    \quad CFP-10 only & --- & --- & 0.797 & \textbf{0.987} & 0.534 & --- \\
    \quad ESAT-6 only & --- & --- & 0.708 & 0.511 & \textbf{0.981} & --- \\
    \bottomrule
  \end{tabular}}
\end{table}

\subsection{Independent Test Set Validation}

On the held-out test set, CatBoost achieved the highest accuracy
(95.6\%) with 93.3\% sensitivity and 98.8\% specificity, while
HistGradientBoosting achieved the highest AUROC of 0.980 (95\% CI
0.967--0.993).
Confusion matrices and ROC curves for all three top models are shown
in Figure~\ref{fig:roc_confusion}.
The DeLong test indicated no statistically significant difference
between the AUROC of CatBoost and HistGradientBoosting
($p$~=~[XX]).
At the operating point that maximizes the Youden index,
CatBoost yielded a positive predictive value (PPV) for TB of
[XX]\% and a negative predictive value (NPV) for TB (i.e., a
correctly identified NTM case) of [XX]\%.

\subsection{Ablation Studies: Mechanistic Basis of Feature
  Complementarity}
\label{sec:ablation}

The ablation experiments reported here are designed as \textit{proofs
of necessity}: each subset is evaluated in isolation not to propose
its clinical use, but to demonstrate that the full 12-feature
combination is required to achieve jointly high sensitivity and
specificity.
Using any single subset in isolation degrades performance
substantially, which in turn confirms that the full feature set
captures complementary, non-redundant diagnostic information.

\paragraph{Charge-state ablation.}
Singly charged (z~=~1) features alone achieved 100\% sensitivity
but only 27.3\% specificity, whereas doubly charged (z~=~2) features
alone maintained 98.8\% specificity but recovered only 46.6\%
sensitivity (Table~\ref{tab:cv_results}).
This asymmetric performance profile is explained by two concurrent
physical and analytical mechanisms.

First, the z~=~1 and z~=~2 ions of the same protein are detected
at $m/z$ values that differ by a factor of two; because the
extraction windows are defined in ppm, the absolute Da width of the
z~=~1 window is approximately twice that of the z~=~2 window
(e.g., \SI{20.2}{\dalton} vs.\ \SI{10.1}{\dalton} for CFP-10).
After TIC normalization, this wider window captures a larger
proportion of diffuse spectral background present in NTM spectra
(matrix clusters, lipid fragments), which can produce non-zero
feature values in the absence of a true CFP-10 or ESAT-6 signal.
The classifier, trained to recognize elevated z~=~1 feature values
as indicative of TB, is consequently misled by these background
contributions in NTM samples, producing the high false-positive rate
(72.7\% of NTM cases misclassified) observed when z~=~1 features
are used in isolation.

Second, z~=~2 ions are secondary ionization products that form
only when the parent protein concentration in the MALDI spot exceeds
a critical threshold.\cite{Gross2004MS}
Their absolute intensities are intrinsically lower than the
corresponding z~=~1 ions, and after TIC normalization they can
fall below the local median background level --- yielding zero-valued
features even when the parent protein is genuinely present at low
but clinically relevant concentrations.
This mechanism accounts for the high false-negative rate (53.4\% of
TB cases missed) when z~=~2 features are used in isolation.

The combination of both charge states is therefore essential:
z~=~1 features ensure that low-concentration TB-specific signals
are not missed (sensitivity driver), while z~=~2 features act as
high-confidence confirmatory signals that suppress false positives
(specificity driver).
Neither charge state is redundant; they occupy complementary
positions in the sensitivity--specificity trade-off space.

\paragraph{Protein-family ablation.}
CFP-10 features alone achieved 98.7\% sensitivity but only 53.4\%
specificity, whereas ESAT-6 features alone achieved 98.1\%
specificity but only 51.1\% sensitivity
(Table~\ref{tab:cv_results}).
This orthogonal performance profile reflects fundamental differences
in the biology and abundance of the two antigen families.
CFP-10 (\textit{Rv3874}) is among the most abundantly secreted
proteins in \textit{M.~tuberculosis} culture filtrate,\cite{Sorensen1995}
producing consistently strong MALDI-TOF signals that are detected
in nearly all TB samples but that, at the extraction threshold
used here, also generate non-zero responses in a proportion of
NTM spectra that contain structurally similar low-molecular-mass
proteins or isobaric contaminants.
ESAT-6 (\textit{Rv3875}), by contrast, is expressed at lower
levels and its secretion is tightly coupled to the ESX-1 type VII
secretion system, which is absent from most NTM
species;\cite{Berthet1998} consequently, a detectable ESAT-6 signal
constitutes strong evidence of \textit{M.~tuberculosis} complex,
but the signal is not present in all TB isolates at the sensitivity
threshold of the extraction method.
The ML model learns to exploit this complementarity: CFP-10 flags
candidate TB cases (high recall), and ESAT-6 confirms or rejects
them (high precision), achieving jointly $>$93\% sensitivity and
$>$98\% specificity in the full-feature configuration.
These findings demonstrate that neither antigen family alone is
sufficient for robust TB/NTM discrimination, and that the diagnostic
value of the pipeline resides precisely in the integration of all
12 targeted features.

\subsection{Feature Importance and Interpretability}

SHAP analysis of the CatBoost model (Figure~\ref{fig:shap}) identified
CFP-10* (\textit{m/z} 10,660) and CFP-10 (\textit{m/z} 10,100) as
the features with the highest mean absolute SHAP values, consistent
with the fold-change analysis (61.4$\times$ and [XX]$\times$,
respectively).
Among ESAT-6 features, [ESAT-6\_1] and [ESAT-6*\_2\_z2] showed the
highest importance in correctly classifying NTM cases,
supporting the interpretation of ESAT-6 as the specificity-driving
antigen family.
The interpretability of the 12-feature model facilitates clinical
communication: a predicted TB case can be associated with the
specific antigen peaks that drove the classification,
enabling clinicians to link the algorithmic decision to a
biologically meaningful signal.

\subsection{Comparison with Standard MALDI-TOF Identification}

% TODO: Add results comparing the ML pipeline against the
% standard Bruker MALDI Biotyper or VITEK MS score approach.
% This section is essential for peer review.
[To be completed: comparison of the proposed pipeline against
standard MALDI-TOF species identification (Biotyper score threshold
$\geq$2.0) and against a simple CFP-10 peak-detection threshold.
This comparison will quantify the diagnostic added value of the
ML layer.]

\subsection{Limitations}

Several limitations of this study should be acknowledged.
First, all spectra were acquired on a single instrument platform
at a single center; inter-instrument reproducibility across different
MALDI-TOF platforms (e.g., Bruker rapifleX, VITEK MS) has not been
empirically validated and warrants prospective multi-site
evaluation.
Second, the hyperparameter search was limited to default
configurations during the initial 43-model screening; further
performance gains may be achievable through systematic optimization
of the top-performing model.
Third, the demographic composition and geographic origin of the cohort
are not yet fully characterized; potential confounders such as
strain diversity (including multidrug-resistant TB), HIV co-infection
status, and prior BCG vaccination may influence spectral profiles
and should be explored in future studies.
Fourth, the model was trained and validated on culture-positive
specimens; its performance on direct clinical samples without culture
amplification is unknown.

%% ═══════════════════════════════════════════════════════════════════
\section{Conclusions}
%% ═══════════════════════════════════════════════════════════════════

We have described and validated a targeted MALDI-TOF MS analytical
pipeline for differential diagnosis of \textit{M.~tuberculosis} from
NTM using 12 proteomic features derived from the CFP-10 and ESAT-6
antigen complex.
The key methodological innovation is the application of the Continuous
Line Segment Algorithm in the $\sqrt{m/z}$ domain for morphological
baseline correction, which eliminates systematic mass-dependent
artefacts that compromise standard ALS-based preprocessing in the
high-mass region.
A patient-stratified benchmark of 43 machine learning algorithms
identified ensemble gradient-boosting models as optimal, with
the best classifier achieving AUROC~=~0.983 and 95.6\% accuracy
on an independent held-out test set.
Ablation experiments revealed that CFP-10 features confer sensitivity
($\geq$99\%) while ESAT-6 features confer specificity ($\geq$98\%),
and that singly and doubly charged ions are complementary rather
than redundant.
Because the pipeline operates exclusively on standard MALDI-TOF
spectra acquired using existing laboratory infrastructure,
with a computational runtime below 5 minutes, it is immediately
deployable in clinical microbiology laboratories and does not
require additional reagents, instrumentation, or specialist
expertise.
Prospective multi-center validation across different MALDI-TOF
platforms and patient populations is warranted to establish the
generalizability of these findings prior to clinical translation.

%% ═══════════════════════════════════════════════════════════════════
\begin{acknowledgement}
% TODO: Add funding sources, grant numbers, acknowledgments.
% Format: "This work was supported by [Funding Agency] (grant number
% [XXXX] to [initials])."
% Do NOT list personal acknowledgments; ACS requires only scientific
% and funding contributions.
[Funding statement to be added.]
\end{acknowledgement}

%% ── Supporting Information ─────────────────────────────────────────
\begin{suppinfo}
The following files are available free of charge:

\textbf{Supporting Information (PDF):}
Table~S1: Complete list of 43 machine learning model configurations
and cross-validation metrics.
Table~S2: Quantitative comparison of CLSA vs.\ ALS baseline
correction SNR improvement across all 12 biomarker windows.
Figure~S1: Representative comparison of CLSA and ALS baseline
correction showing systematic underestimation by ALS at
\textgreater\SI{8000}{\dalton}.
Figure~S2: Confusion matrices for all three top-ranked classifiers
on the independent test set.
Figure~S3: SHAP beeswarm plot for the full CatBoost model.
\end{suppinfo}

%% ── Bibliography ────────────────────────────────────────────────────
\bibliography{references}

%% ═══════════════════════════════════════════════════════════════════
%% FIGURE CAPTIONS
%% (In achemso the captions are placed within float environments;
%%  they are listed here for editorial reference only.)
%% ═══════════════════════════════════════════════════════════════════

%% Figure 1 — Mean spectra comparison
\begin{figure}
  \centering
  \includegraphics[width=\linewidth]{figures/fig1_mean_spectra}
  \caption{Mean MALDI-TOF MS spectra ($\pm$1 SD) for TB (\textcolor{blue}{blue})
    and NTM (\textcolor[RGB]{230,159,0}{orange}) isolates after preprocessing,
    over the \SIrange{3}{15}{\kilo\dalton} range.
    Insets show expanded views of the (a)~CFP-10/CFP-10* region
    (\SIrange{9.8}{11.0}{\kilo\dalton}) and (b)~ESAT-6/ESAT-6* region
    (\SIrange{7.7}{10.0}{\kilo\dalton}).
    Dashed vertical lines indicate the theoretical $m/z$ of all 12
    targeted biomarker features (Table~\ref{tab:biomarkers}).}
  \label{fig:mean_spectra}
\end{figure}

%% Figure 2 — Preprocessing pipeline
\begin{figure}
  \centering
  \includegraphics[width=\linewidth]{figures/fig2_preprocessing}
  \caption{Sequential preprocessing of a representative TB MALDI-TOF
    spectrum.
    (a)~Raw spectrum; (b)~after Savitzky--Golay smoothing (window~=~21,
    order~=~3); (c)~CLSA baseline estimate (red) overlaid on smoothed
    spectrum; (d)~baseline-subtracted and TIC-normalized spectrum with
    biomarker extraction windows highlighted.
    The $\sqrt{m/z}$ transformation applied within the CLSA algorithm
    normalizes peak widths across the mass range, preventing
    under-correction at higher masses.}
  \label{fig:preprocessing}
\end{figure}

%% Figure 3 — Biomarker feature distributions
\begin{figure}
  \centering
  \includegraphics[width=\linewidth]{figures/fig3_boxplots}
  \caption{Distribution of the 12 biomarker feature intensities
    (after TIC normalization and local background subtraction) for TB
    (blue) and NTM (orange) classes.
    Boxes indicate median and interquartile range; whiskers extend to
    1.5 $\times$ IQR.
    Asterisks indicate statistical significance after Bonferroni
    correction ($^{*}p<0.05$; $^{**}p<0.01$; $^{***}p<0.001$;
    Mann--Whitney U test).
    Fold changes (TB/NTM mean) are reported above each panel.}
  \label{fig:boxplots}
\end{figure}

%% Figure 4 — PCA
\begin{figure}
  \centering
  \includegraphics[width=0.6\linewidth]{figures/fig4_pca}
  \caption{Principal component analysis (PCA) of the 12-biomarker
    feature matrix.
    Each point represents one spectrum, colored by diagnostic class
    (TB: blue; NTM: orange).
    Ellipses represent the 95\% confidence region assuming
    bivariate normality.
    Percentages indicate variance explained by PC1 and PC2.}
  \label{fig:pca}
\end{figure}

%% Figure 5 — ML benchmark
\begin{figure}
  \centering
  \includegraphics[width=\linewidth]{figures/fig5_benchmark}
  \caption{Cross-validation AUROC for all 43 machine learning
    configurations, grouped by algorithm family.
    Points indicate the mean of the 5 folds; error bars indicate
    $\pm$1 SD.
    Colors distinguish algorithm families as indicated in the legend.
    The dashed horizontal line marks the performance of the
    stratified dummy classifier (no-information baseline).}
  \label{fig:benchmark}
\end{figure}

%% Figure 6 — ROC curves and confusion matrices
\begin{figure}
  \centering
  \includegraphics[width=\linewidth]{figures/fig6_roc_confusion}
  \caption{Independent test set performance of the three top-ranked
    classifiers.
    (a--c)~Confusion matrices (rows: true label; columns: predicted
    label); values are counts and row-normalized percentages.
    (d)~Receiver operating characteristic (ROC) curves with 95\%
    confidence bands computed by DeLong bootstrapping.
    Diagonal dashed line: no-discrimination reference.}
  \label{fig:roc_confusion}
\end{figure}

%% Figure 7 — SHAP feature importance
\begin{figure}
  \centering
  \includegraphics[width=0.7\linewidth]{figures/fig7_shap}
  \caption{SHAP beeswarm plot for the CatBoost classifier on the
    independent test set.
    Each point represents one spectrum; the $x$-axis shows the SHAP
    value (contribution to the log-odds of a TB prediction);
    color encodes the feature value (blue: low; red: high).
    Features are ordered by mean absolute SHAP value (top = most
    important).}
  \label{fig:shap}
\end{figure}

\end{document}
